% Chapter 11: DevOps Pipeline and Release Engineering
\chapter{DEVOPS PIPELINE AND RELEASE ENGINEERING}

Nebula is not only a software product but also a deployment system. For that reason, the surrounding DevOps pipeline is part of the technical contribution of the project. A control plane that manages infrastructure across clouds must itself be delivered through a repeatable, auditable, and low-risk release process. This chapter describes the release engineering practices that support that goal.

\section{Source Control and Branching Strategy}
The codebase is best maintained through a branch strategy that separates active development from release-ready code. A simple and effective structure includes:
\begin{itemize}
    \item \textbf{main}: production-ready branch with protected merges,
    \item \textbf{develop}: integration branch for completed features,
    \item \textbf{feature/*}: isolated task branches,
    \item \textbf{hotfix/*}: urgent corrections to production-impacting issues.
\end{itemize}

This model reduces release uncertainty because infrastructure changes, frontend changes, and backend changes can be reviewed together before they reach a shared environment. It also supports traceability: every release can be tied back to a known set of commits, pull requests, and approvals.

\section{Build Artifacts and Versioning}
Nebula produces several deployable artifacts:
\begin{enumerate}
    \item frontend static assets,
    \item backend API container image,
    \item worker container image,
    \item Terraform module templates and shared configuration bundles,
    \item database migration scripts.
\end{enumerate}

All artifacts should be versioned consistently using a release tag such as \texttt{v1.4.0}. The critical idea is that the API image, worker image, and schema migration state should correspond to the same release identity. Without that discipline, a debugging session becomes unnecessarily difficult because different parts of the system may be running incompatible assumptions.

\section{Continuous Integration Stages}
The continuous integration pipeline acts as the first quality gate. At minimum, it should perform:
\begin{itemize}
    \item backend static analysis and unit testing,
    \item frontend linting and build verification,
    \item Terraform syntax validation for bundled templates,
    \item schema migration validation,
    \item API contract or smoke-style integration tests.
\end{itemize}

An important design decision is to treat Terraform templates as first-class code artifacts. Even if the backend and frontend pass their tests, a malformed module template can still break the actual business function of the platform. Therefore, the IaC layer should be validated alongside the application layer.

\subsection{Illustrative Pipeline Definition}
\begin{lstlisting}[language=Python, caption=Representative CI Pipeline Skeleton]
stages:
  - validate
  - test
  - package
  - deploy

validate_backend:
  stage: validate
  script:
    - python -m pytest tests/unit

validate_frontend:
  stage: validate
  script:
    - npm ci
    - npm run build

validate_iac:
  stage: validate
  script:
    - terraform fmt -check
    - terraform validate

package_api:
  stage: package
  script:
    - docker build -t nebula-api:$CI_COMMIT_TAG backend

package_worker:
  stage: package
  script:
    - docker build -t nebula-worker:$CI_COMMIT_TAG backend
\end{lstlisting}

\section{Environment Promotion Workflow}
Releases should move progressively from lower-risk to higher-risk environments. A common pattern is:
\begin{enumerate}
    \item merge reviewed code into the integration branch,
    \item deploy automatically to a shared test environment,
    \item execute smoke tests for login, queueing, and provider connectivity,
    \item promote the same versioned artifact to staging,
    \item require approval for production deployment,
    \item monitor early production metrics closely after rollout.
\end{enumerate}

This progression matters because a multi-cloud control plane can fail in subtle ways that ordinary unit tests do not catch. A release may look healthy until it attempts provider authentication, state reconciliation, or asynchronous task execution with live infrastructure APIs. Promotion through increasingly realistic environments is therefore a core risk-reduction strategy.

\section{Secrets and Configuration Management}
The release pipeline must distinguish between code and environment-specific secret material. The following items should never be embedded in source control:
\begin{itemize}
    \item master encryption keys,
    \item production database credentials,
    \item provider service principal secrets,
    \item external AI or API tokens.
\end{itemize}

Configuration should instead be injected at runtime through environment variables, secret stores, or orchestrator-native secret mechanisms. This approach supports safer rotation, reduces accidental leakage, and keeps the codebase portable between environments.

\section{Containerization Strategy}
Containerization is valuable for Nebula because it stabilizes execution across environments. The API, worker, database, and broker can be composed predictably, making it easier for developers to reproduce production-adjacent behavior on local systems. The main benefits are:
\begin{itemize}
    \item consistent runtime dependencies,
    \item simpler onboarding,
    \item versioned deployment artifacts,
    \item clear service boundaries,
    \item easier rollback to known-good images.
\end{itemize}

However, containerization alone does not guarantee operational quality. Images still need vulnerability review, startup probes, health endpoints, and logging standards. This is especially relevant for the worker container, which handles secret decryption and Terraform execution.

\section{Database Migration Discipline}
Schema changes are one of the most common sources of release fragility. Nebula stores users, credentials, task records, and resource state, so a careless migration can affect both security and availability. A responsible migration policy includes:
\begin{enumerate}
    \item forward-only migrations where possible,
    \item explicit review of nullable-to-non-null changes,
    \item data backfills as controlled release steps,
    \item backups before destructive changes,
    \item smoke tests against the migrated schema.
\end{enumerate}

\section{Rollback Strategy}
Rollback must be defined differently for code, configuration, and infrastructure:
\begin{itemize}
    \item \textbf{Code rollback}: redeploy previous API and worker images.
    \item \textbf{Configuration rollback}: restore known-good environment settings and feature flags.
    \item \textbf{Infrastructure rollback}: apply carefully reviewed reverse changes rather than assuming all cloud resources can be reverted instantly.
\end{itemize}
This distinction is important because a release may fail at one layer but remain healthy at another. For example, the API code may be safe to roll back while the database schema remains forward compatible.

\section{Release Readiness Checklist}
\begin{table}[H]
\centering
\caption{Release Gate Criteria}
\label{tab:release_gate}
\begin{tabular}{|p{4.0cm}|p{2.2cm}|p{5.0cm}|}
\hline
\textbf{Check} & \textbf{Required} & \textbf{Rationale} \\ \hline
Unit and integration tests pass & Yes & baseline quality signal \\ \hline
Terraform templates validate & Yes & prevents broken provisioning flows \\ \hline
Migration reviewed and backed up & Yes & protects persistent state \\ \hline
Secrets resolved in target environment & Yes & avoids startup and runtime failures \\ \hline
Staging smoke test completed & Yes & confirms end-to-end functionality \\ \hline
Operational owner assigned & Yes & ensures release accountability \\ \hline
Rollback note prepared & Recommended & accelerates incident response \\ \hline
\end{tabular}
\end{table}

\section{Discussion}
The DevOps pipeline is not separate from the product; it is part of the product's reliability envelope. Nebula aims to reduce operational complexity for its users, and it must apply the same discipline to itself. By versioning artifacts, validating infrastructure templates, managing secrets outside source code, and promoting releases through controlled environments, the platform becomes easier to trust and easier to maintain. These practices also strengthen the academic value of the project because they show awareness of the full software lifecycle rather than implementation in isolation.

\newpage
